\documentclass{amsart}
% For dealing with references we use the comment environment
\usepackage{verbatim}
\newenvironment{reference}{\comment}{\endcomment}
%\newenvironment{reference}{}{}
\newenvironment{slogan}{\comment}{\endcomment}
\newenvironment{history}{\comment}{\endcomment}

% For commutative diagrams we use Xy-pic
\usepackage[all]{xy}

% We use 2cell for 2-commutative diagrams.
\xyoption{2cell}
\UseAllTwocells

% We use multicol for the list of chapters between chapters
\usepackage{multicol}

% This is generally recommended for better output
\usepackage{lmodern}
\usepackage[T1]{fontenc}

% For cross-file-references

% Package for hypertext links:
\usepackage{hyperref}

% For any local file, say "hello.tex" you want to link to please

% Theorem environments.
%
\theoremstyle{plain}
\newtheorem{theorem}[subsection]{Theorem}
\newtheorem{proposition}[subsection]{Proposition}
\newtheorem{lemma}[subsection]{Lemma}

\theoremstyle{definition}
\newtheorem{definition}[subsection]{Definition}
\newtheorem{example}[subsection]{Example}
\newtheorem{exercise}[subsection]{Exercise}
\newtheorem{situation}[subsection]{Situation}

\theoremstyle{remark}
\newtheorem{remark}[subsection]{Remark}
\newtheorem{remarks}[subsection]{Remarks}

\numberwithin{equation}{subsection}

% Macros
%
\def\lim{\mathop{\mathrm{lim}}\nolimits}
\def\colim{\mathop{\mathrm{colim}}\nolimits}
\def\Spec{\mathop{\mathrm{Spec}}}
\def\Hom{\mathop{\mathrm{Hom}}\nolimits}
\def\Ext{\mathop{\mathrm{Ext}}\nolimits}
\def\SheafHom{\mathop{\mathcal{H}\!\mathit{om}}\nolimits}
\def\SheafExt{\mathop{\mathcal{E}\!\mathit{xt}}\nolimits}
\def\Sch{\mathit{Sch}}
\def\Mor{\mathop{\mathrm{Mor}}\nolimits}
\def\Ob{\mathop{\mathrm{Ob}}\nolimits}
\def\Sh{\mathop{\mathit{Sh}}\nolimits}
\def\NL{\mathop{N\!L}\nolimits}
\def\CH{\mathop{\mathrm{CH}}\nolimits}
\def\proetale{{pro\text{-}\acute{e}tale}}
\def\etale{{\acute{e}tale}}
\def\QCoh{\mathit{QCoh}}
\def\Ker{\mathop{\mathrm{Ker}}}
\def\Im{\mathop{\mathrm{Im}}}
\def\Coker{\mathop{\mathrm{Coker}}}
\def\Coim{\mathop{\mathrm{Coim}}}

% Boxtimes
%
\DeclareMathSymbol{\boxtimes}{\mathbin}{AMSa}{"02}

%
% Macros for moduli stacks/spaces
%
\def\QCohstack{\mathcal{QC}\!\mathit{oh}}
\def\Cohstack{\mathcal{C}\!\mathit{oh}}
\def\Spacesstack{\mathcal{S}\!\mathit{paces}}
\def\Quotfunctor{\mathrm{Quot}}
\def\Hilbfunctor{\mathrm{Hilb}}
\def\Curvesstack{\mathcal{C}\!\mathit{urves}}
\def\Polarizedstack{\mathcal{P}\!\mathit{olarized}}
\def\Complexesstack{\mathcal{C}\!\mathit{omplexes}}
% \Pic is the operator that assigns to X its picard group, usage \Pic(X)
% \Picardstack_{X/B} denotes the Picard stack of X over B
% \Picardfunctor_{X/B} denotes the Picard functor of X over B
\def\Pic{\mathop{\mathrm{Pic}}\nolimits}
\def\Picardstack{\mathcal{P}\!\mathit{ic}}
\def\Picardfunctor{\mathrm{Pic}}
\def\Deformationcategory{\mathcal{D}\!\mathit{ef}}

\setlength{\emergencystretch}{3em}
\begin{document}
\title[Stacks Project, Tag 05TW]{360 --- Existence of strict filtered-injective resolutions of bounded-below filtered complexes}
\maketitle
\noindent Copyright (C) 2005--2025 Johan de Jong. Stacks Project authors. Source: \href{https://stacks.math.columbia.edu/tag/05TW}{Tag 05TW}.

\noindent Editorial difficulty note (not author text). Difficulty H2: The proof constructs compatible resolutions of each kernel and coimage, combines them through split exact sequences, and builds a double complex with square-zero horizontal differential. It then proves its totalization is a filtered quasi-isomorphism by comparing associated graded total complexes. Resolving the exact filtered category requires these additional compatibility constructions..

\noindent Editorial provenance note (not author text). History: excerpt from revision \texttt{a0444\allowbreak 6e57e\allowbreak c1fbc\allowbreak 252a8\allowbreak 71afc\allowbreak ec775\allowbreak 2fb28\allowbreak 07b14}, derived.tex, lines 8095--8194. Reference presentation was adapted; original-author context and core support blocks were retained or added. The mathematical wording is unchanged. Licensed under GNU FDL 1.2 or later; no invariant sections or cover texts. See ../licenses/COPYING. Context blocks reproduce author definitions and supporting results; they are not additional counted problems.

\begin{definition}
\label{definition-filtered-complexes-notation}
Let $\mathcal{A}$ be an abelian category.
We say an object $I$ of $\text{Fil}^f(\mathcal{A})$
is {\it filtered injective} if each $\text{gr}^p(I)$ is
an injective object of $\mathcal{A}$.
\end{definition}

\section{Filtered derived category and injective resolutions}
\label{section-filtered-derived}

\noindent
Let $\mathcal{A}$ be an abelian category. In this section we will show
that if $\mathcal{A}$ has enough injectives, then so does
the category $\text{Fil}^f(\mathcal{A})$ in some sense. One
can use this observation to compute in the filtered derived category
of $\mathcal{A}$.

\medskip\noindent
The category $\text{Fil}^f(\mathcal{A})$ is an example of an
exact category, see
Injectives, Remark \href{https://stacks.math.columbia.edu/tag/05SF}{remark embed exact category (Tag 05SF)}.
A special role is played by the strict morphisms, see
Homology, Definition \ref{homology-definition-strict},
i.e., the morphisms $f$ such that $\Coim(f) = \Im(f)$.
We will say that a complex $A \to B \to C$ in $\text{Fil}^f(\mathcal{A})$ is
{\it exact} if the sequence $\text{gr}(A) \to \text{gr}(B) \to \text{gr}(C)$
is exact in $\mathcal{A}$. This implies that $A \to B$ and $B \to C$
are strict morphisms, see
Homology, Lemma \href{https://stacks.math.columbia.edu/tag/05QH}{lemma filtered acyclic (Tag 05QH)}.



\begin{definition}
\label{homology-definition-strict}
Let $\mathcal{A}$ be an abelian category.
A morphism $f : A \to B$ of filtered objects of $\mathcal{A}$ is
said to be {\it strict} if $f(F^iA) = f(A) \cap F^iB$ for
all $i \in \mathbf{Z}$.
\end{definition}

\begin{definition}
\label{homology-definition-associated-simple-complex}
Let $\mathcal{A}$ be an additive category.
Let $A^{\bullet, \bullet}$ be a double complex.
The {\it associated simple complex}, denoted $sA^\bullet$, also
often called the {\it associated total complex}, denoted
$\text{Tot}(A^{\bullet, \bullet})$, is
given by
$$
sA^n = \text{Tot}^n(A^{\bullet, \bullet}) =
\bigoplus\nolimits_{n = p + q} A^{p, q}
$$
(if it exists) with differential
$$
d_{sA^\bullet}^n = d_{\text{Tot}(A^{\bullet, \bullet})}^n =
\sum\nolimits_{n = p + q} (d_1^{p, q} + (-1)^p d_2^{p, q})
$$
\end{definition}

\begin{lemma}
\label{lemma-right-resolution-by-filtered-injectives}
Let $\mathcal{A}$ be an abelian category with enough injectives.
For every $K^\bullet \in K^{+}(\text{Fil}^f(\mathcal{A}))$
there exists a filtered quasi-isomorphism $K^\bullet \to I^\bullet$
with $I^\bullet$ bounded below,
each $I^n$ a filtered injective object, and
each $K^n \to I^n$ a strict monomorphism.
\end{lemma}

\begin{proof}
After replacing $K^\bullet$ by a shift (which is harmless for the proof)
we may assume that $K^n = 0$ for $n < 0$. Consider the
short exact sequences
$$
\begin{matrix}
0 \to \Ker(d_K^0) \to K^0 \to \Coim(d_K^0) \to 0 \\
0 \to \Ker(d_K^1) \to K^1 \to \Coim(d_K^1) \to 0 \\
0 \to \Ker(d_K^2) \to K^2 \to \Coim(d_K^2) \to 0 \\
\ldots
\end{matrix}
$$
of the exact category $\text{Fil}^f(\mathcal{A})$
and the maps $u_i : \Coim(d_K^i) \to \Ker(d_K^{i + 1})$.
For each $i \geq 0$ we may choose filtered quasi-isomorphisms
$$
\begin{matrix}
\Ker(d_K^i)[0] \to I_{ker, i}^\bullet \\
\Coim(d_K^i)[0] \to I_{coim, i}^\bullet
\end{matrix}
$$
with $I_{ker, i}^n, I_{coim, i}^n$ filtered injective and zero for $n < 0$, see
Lemma \href{https://stacks.math.columbia.edu/tag/05TT}{lemma filtered injective right resolution single object (Tag 05TT)}.
By
Lemma \href{https://stacks.math.columbia.edu/tag/05TU}{lemma filtered injective right resolution map (Tag 05TU)}
we may lift $u_i$ to a morphism of complexes
$u_i^\bullet : I_{coim, i}^\bullet \to I_{ker, i + 1}^\bullet$.
Finally, for each $i \geq 0$ we may complete the diagrams
$$
\xymatrix{
0 \ar[r] &
\Ker(d_K^i)[0] \ar[r] \ar[d] &
K^i[0] \ar[r] \ar[d] &
\Coim(d_K^i)[0] \ar[r] \ar[d] &
0 \\
0 \ar[r] &
I_{ker, i}^\bullet \ar[r]^{\alpha_i} &
I_i^\bullet \ar[r]^{\beta_i} &
I_{coim, i}^\bullet \ar[r] &
0
}
$$
with the lower sequence a termwise split exact sequence, see
Lemma \href{https://stacks.math.columbia.edu/tag/05TV}{lemma filtered injective right resolution ses (Tag 05TV)}.
For $i \geq 0$ set $d_i : I_i^\bullet \to I_{i + 1}^\bullet$
equal to $d_i =  \alpha_{i + 1} \circ u_i^\bullet \circ \beta_i$.
Note that $d_i \circ d_{i - 1} = 0$ because
$\beta_i \circ \alpha_i = 0$. Hence we have constructed
a commutative diagram
$$
\xymatrix{
I_0^\bullet \ar[r] &
I_1^\bullet \ar[r] &
I_2^\bullet \ar[r] & \ldots \\
K^0[0] \ar[r] \ar[u] &
K^1[0] \ar[r] \ar[u] &
K^2[0] \ar[r] \ar[u] &
\ldots
}
$$
Here the vertical arrows are filtered quasi-isomorphisms.
The upper row is a complex of complexes and each complex consists of
filtered injective objects with no nonzero objects in degree $< 0$.
Thus we obtain a double complex by setting $I^{a, b} = I_a^b$ and using
$$
d_1^{a, b} : I^{a, b} = I_a^b \to I_{a + 1}^b = I^{a + 1, b}
$$
the map $d_a^b$ and using for
$$
d_2^{a, b} : I^{a, b} = I_a^b \to I_a^{b + 1} = I^{a, b + 1}
$$
the map $d_{I_a}^b$. Denote $\text{Tot}(I^{\bullet, \bullet})$
the total complex associated to this double complex, see
Homology, Definition \ref{homology-definition-associated-simple-complex}.
Observe that the maps $K^n[0] \to I_n^\bullet$ come from maps
$K^n \to I^{n, 0}$ which give rise to a map of complexes
$$
K^\bullet \longrightarrow \text{Tot}(I^{\bullet, \bullet})
$$
We claim this is a filtered quasi-isomorphism.
As $\text{gr}(-)$ is an additive functor, we see that
$\text{gr}(\text{Tot}(I^{\bullet, \bullet})) =
\text{Tot}(\text{gr}(I^{\bullet, \bullet}))$.
Thus we can use
Homology,
Lemma \href{https://stacks.math.columbia.edu/tag/0133}{lemma double complex gives resolution (Tag 0133)}
to conclude that
$\text{gr}(K^\bullet) \to \text{gr}(\text{Tot}(I^{\bullet, \bullet}))$
is a quasi-isomorphism as desired.
\end{proof}
\end{document}
